% Generated draft: requires mathematical and source review.
\documentclass[11pt]{article}
\usepackage[margin=25mm]{geometry}
\usepackage{fontspec}
\setmainfont{Latin Modern Roman}
\usepackage{amsmath,amssymb,mathtools,unicode-math}
\setmathfont{Latin Modern Math}
\usepackage{enumitem,needspace,xurl,hyperref}
\setlength{\parindent}{0pt}
\setlength{\parskip}{5pt}
\setlength{\emergencystretch}{3em}
\newcommand{\N}{\mathbb{N}}
\newcommand{\Z}{\mathbb{Z}}
\newcommand{\Q}{\mathbb{Q}}
\newcommand{\R}{\mathbb{R}}
\newcommand{\C}{\mathbb{C}}
\newcommand{\F}{\mathbb{F}}
\newcommand{\A}{\mathbb{A}}
\newcommand{\PP}{\mathbb{P}}
\newcommand{\E}{\mathbb{E}}
\newcommand{\eps}{\varepsilon}
\newcommand{\dd}{\,\mathrm{d}}
\begin{document}
\section*{UnsolvedMath ID: 30002541}
\textbf{Source problem code:} catalogue entry
\subsection{Definitions and mathematical statement}
Let \(\mathcal G\) denote a class of finite graphs. A block is a maximal connected subgraph having no cut vertex; under the usual convention, bridges give blocks isomorphic to \(K_2\), and isolated vertices may be treated separately. In the usual meaning of \emph{block-stable}, membership of a connected graph in the class is determined by membership of its blocks in a specified allowed collection. The supplied record does not specify the conventions concerning isolated vertices, connected versus arbitrary graphs, or closure under taking disjoint unions.

For a finite connected graph \(G\), write \(V(G)\) for its vertex set and define the graph distance by
\[
 d_G(u,v)=\min\{\text{number of edges in a path from }u\text{ to }v\},\qquad u,v\in V(G).
\]
A distance normalization is a deterministic positive sequence \((a_n)\); it gives the rescaled metric space
\[
 \bigl(V(G_n),a_n d_{G_n}\bigr).
\]
Here \(G_n\) would be a random graph of order \(n\) in the relevant class, for those orders for which the chosen probability model exists. Its probability law is not specified by the source. If a measured scaling limit is intended, one must additionally specify a vertex measure and a topology for convergence.

\textbf{Problem as supplied.} Determine the asymptotic behavior of graph distances and the scaling limits for random graphs drawn from block-stable classes that are not subcritical, with random planar graphs as the principal case. More explicitly, after specifying the class and random graph model, determine appropriate distance normalizations \(a_n\), establish whether the resulting random metric spaces have scaling limits in a specified topology, and identify such limits where they exist.

\textbf{Unresolved formulation qualifications.} The source supplies neither the analytic definition of subcriticality nor the hypotheses on the graph classes under consideration. Thus ``not subcritical'' is retained as a source qualification, rather than replaced by an invented analytic condition. No universal normalization or common limiting space for all such classes is asserted. The source also leaves unspecified whether distances mean distances between sampled vertices, the diameter, or the full random metric space; whether planar graphs are connected, simple, labelled, or unlabelled; and whether sampling is uniform or weighted. For disconnected graphs, ordinary graph distance does not give a finite metric on the entire vertex set, so a component convention or another distance convention is required. Consequently, this is a research target requiring further hypotheses, rather than a fully specified convergence conjecture.
\subsection{Short English statement}
Determine how distances scale, and what limiting random metric spaces arise, for random graphs in block-stable classes outside the subcritical regime, especially planar graphs. The source leaves the sampling model, connectivity assumptions, meaning of subcriticality, distance observable, and convergence topology unspecified. These choices require human mathematical review before a precise convergence claim can be formulated.
\subsection{Sources}
\begin{itemize}
\item Supplied local UnsolvedMath registry record, permanent ID \(30002541\), title \emph{Scaling Limits of Nonsubcritical Block Stable Random Graphs}; exact source problem code \texttt{catalogue entry}. Collection: Oberwolfach Reports Open Problems. The record gives proposed year 2014 and names no individual proposer.
\item Supplied report citation: \emph{Enumerative Combinatorics} (2014), Oberwolfach Reports, DOI 10.4171/owr/2014/12, pp.\ 635--720. \url{https://doi.org/10.4171/owr/2014/12}. The supplied background describes these reports as documenting mathematical workshops and their open-problem discussions; no exact problem page or verbatim report passage is supplied.
\item Supplied literature source, labelled ``Stufler (2024)'': \url{https://onlinelibrary.wiley.com/doi/10.1112/jlms.70018}.
\item Supplied literature source, labelled ``Cubic planar graphs'': \url{https://projecteuclid.org/journals/electronic-journal-of-probability/volume-28/issue-none/Random-cubic-planar-graphs-converge-to-the-Brownian-sphere/10.1214/23-EJP912.pdf}.
\item Supplied literature source, labelled ``Original Oberwolfach report'': \url{https://ems.press/content/serial-article-files/46502?nt=1}.
\item Literature assessment dated exactly \texttt{2026-08-21}: the supplied curation labels the problem open and states that Brownian-map scaling limits are known for important subclasses such as cubic planar graphs, while the extracted general planar or block-stable non-subcritical problem remains unresolved. This is supplied dated curation, not an independently verified assessment. The supplied background explicitly requests independent verification of time-sensitive status and cited claims.
\item The registry labels the statement \texttt{corrected\_verified} and states that its mathematical content was unchanged while notation, typography, and wording were normalized into a self-contained publication form. This is supplied registry metadata, not verification performed for this draft.
\end{itemize}
\end{document}
